\subsubsection{Weighted-factorization control of $\widetilde\Gamma_t^N$}
\label{subsubsec:weighted-factorization-tilde-gamma}

We bound the full jump covariance residual directly in a Hilbert space of
weighted features. The proof has three steps: factor the Hessian kernel,
concentrate the feature covariance in Hilbert--Schmidt norm, and convert
that estimate to the nuclear norm of the finite jump covariance matrix.

Throughout this subsection, fix a population realization in $E_{\rm pop}$
from \cref{lem:population-moment-event}, and write
\[
 \mathbb E_{\mathcal G}[Z]:=\mathbb E[Z\mid\mathcal G_{\rm pop}],
 \qquad
 \mathbb P_{\mathcal G}(A):=\mathbb P(A\mid\mathcal G_{\rm pop}).
\]
The conditional product law in \cref{lemma-gp-jump-implies-conditional-iid}
and the conditional moment bound \eqref{eq:population-uniform-conditional-moments}
give, at every fixed time and its relevant trace,
\begin{equation}
\label{eq:weighted-feature-conditional-input}
 \operatorname{Law}((Y_t^1,\ldots,Y_t^N)\mid\mathcal G_{\rm pop})
 =\mu_t^{\otimes N},\qquad
 (\mathbb E_{\mathcal G}\|Y_t^1\|_{\mathbb R^d}^r)^{1/r}
 \le K_{\rm pop}r^{1/\alpha}\quad(r\ge1).
\end{equation}
Here $K_{\rm pop}$ is the deterministic envelope chosen for this population
event; it may depend on the horizon and its probability budget. We never
condition further on the sample localization event.

For the fixed deterministic radius $D\ge1$, recall
$\mathrm L_D=\mathrm L_{K_{\rm pop},D}$ from
\cref{def:localized-envelope}. The population moment belongs to
$\mathcal V_{K_{\rm pop},D}$ on $E_{\rm pop}$ by
\cref{lemma-reachable-moment-tube}. Thus the moment-Hessian bound
\eqref{eq:local-moment-derivative-bound} is available here without a bound
on the realized sample positions. We retain $\mathrm L_D^2$ explicitly.

\paragraph{Weighted features and the integral Hessian kernel.}
Recall $a_m=e^{-c_am^{\kappa_{\rm dec}}}$ and
$\|h\|_{\mathsf H_a}^2=\sum_ma_m^2\|h_m\|_{\mathsf H_m}^2$ from
\cref{assumption-structured-moment-smoothness}. Also recall the feature map
$J_0:\mathbb R^d\to\mathsf H_a$, $J_0(x)=(\phi_m(x))_{m\ge1}$, from
\cref{lem:localized-weighted-feature-envelopes}, with
$\mathrm D J_0(x)b=(\mathrm D\phi_m(x)b)_{m\ge1}$. Define
\[
 \mathsf E:=\bigoplus_{m\ge1}\mathsf H_m,\qquad
 \mathsf U_a:\mathsf H_a\longrightarrow\mathsf E,
 \quad(\mathsf U_a h)_m=a_mh_m.
\]
The direct sum has its unweighted Hilbert norm, so
$\|\mathsf U_a h\|_{\mathsf E}=\|h\|_{\mathsf H_a}$, and
$(\mathsf U_a^{-1}u)_m=a_m^{-1}u_m$ satisfies
$\|\mathsf U_a^{-1}u\|_{\mathsf H_a}^2=\sum_m\|u_m\|_{\mathsf H_m}^2$.
Thus $\mathsf U_a$ is an isometric isomorphism. Express the derivative
of $J_0$ in these unweighted coordinates by defining
\begin{equation}
\label{eq:weighted-feature-j0-identification}
\begin{aligned}
 \Psi_\infty(x)&:=\mathsf U_a\circ\mathrm D J_0(x)
       :\mathbb R^d\to\mathsf E,\\
 \Psi_\infty(x)b&=\mathsf U_a((\mathrm D\phi_m(x)b)_m)
       =(a_m\mathrm D\phi_m(x)b)_m.
\end{aligned}
\end{equation}
This map is Hilbert--Schmidt:
\begin{equation}
\label{eq:full-weighted-feature-norm}
 \|\Psi_\infty(x)\|_{\mathrm{HS};\,\mathbb R^d\to\mathsf E}^2
 =\sum_{m\ge1}a_m^2\|\mathrm D\phi_m(x)\|_{\mathrm{HS};\,\mathbb R^d\to\mathsf H_m}^2
 \le d\sum_{m\ge1}a_m^2L_m^2(1+\|x\|_{\mathbb R^d}^{m-1})^2<\infty.
\end{equation}
The factor $d$ follows by summing the squared operator bound over an
orthonormal basis of $\mathbb R^d$. Finiteness follows by squaring \cref{lem:weighted-series-estimate},
\eqref{eq:weighted-series-summand-bound}, with $K=1$ and
$D=1+\|x\|_{\mathbb R^d}$, and summing.

The following factorization uses the full feature space $\mathsf E$.

\begin{lemma}[Weighted Hessian factorization]
\label{lemma-structured-weighted-hessian-factorization}
For $M(\mu)\in\mathcal V_{K_{\rm pop},D}$, there is a bounded self-adjoint
operator $\overline A_\infty(\mu)$ on $\mathsf E$ such that
\begin{align}
 \langle u,\overline A_\infty(\mu)v\rangle
 &=\mathrm D^2G(M(\mu))[\mathsf U_a^{-1}u,\mathsf U_a^{-1}v],
 \label{eq:full-weighted-moment-hessian}\\
 \nabla_W^2F(\mu;x,z)
 &=\Psi_\infty(x)^*\overline A_\infty(\mu)\Psi_\infty(z),
 \qquad \|\overline A_\infty(\mu)\|_{\mathsf E\to\mathsf E}\le\mathrm L_D.
 \label{eq:full-weighted-hessian-factorization}
\end{align}
\end{lemma}
\begin{proof}
\noindent\textbf{Step 1: represent the bounded moment bilinear form.}
By the isometry $\mathsf U_a$ and \eqref{eq:local-moment-derivative-bound},
\[
 |\mathrm D^2G(M(\mu))[\mathsf U_a^{-1}u,\mathsf U_a^{-1}v]|
 \le\mathrm L_D\|u\|_{\mathsf E}\|v\|_{\mathsf E}.
\]
For fixed $v\in\mathsf E$, the map
$u\mapsto\mathrm D^2G(M(\mu))[\mathsf U_a^{-1}u,\mathsf U_a^{-1}v]$
is a continuous linear functional of norm at most
$\mathrm L_D\|v\|_{\mathsf E}$. Riesz representation gives its unique representing
vector $\overline A_\infty(\mu)v$. Uniqueness and linearity in $v$
make this a linear operator, with
\[
 \|\overline A_\infty(\mu)v\|_{\mathsf E}
 =\sup_{\|u\|_{\mathsf E}=1}
   |\mathrm D^2G(M(\mu))[\mathsf U_a^{-1}u,\mathsf U_a^{-1}v]|
 \le\mathrm L_D\|v\|_{\mathsf E}.
\]
Symmetry of the second derivative gives
$\langle u,\overline A_\infty(\mu)v\rangle
=\langle\overline A_\infty(\mu)u,v\rangle$,
so this bounded operator is self-adjoint.
\par\noindent\textbf{Step 2: identify the spatial Hessian kernel.}
For $b,c\in\mathbb R^d$, \cref{eq:weighted-feature-j0-identification}
gives $\mathsf U_a^{-1}\Psi_\infty(x)b=\mathrm D J_0(x)b$ and the
corresponding identity at $z$. Substituting these identities into
\cref{eq:weighted-integral-hessian-representation} and using
\cref{eq:full-weighted-moment-hessian} gives
\[
\begin{aligned}
 \langle b,\nabla_W^2F(\mu;x,z)c\rangle_{\mathbb R^d}
 &=\mathrm D^2G(M(\mu))[\mathrm D J_0(x)b,\mathrm D J_0(z)c]\\
 &=\mathrm D^2G(M(\mu))[
       \mathsf U_a^{-1}\Psi_\infty(x)b,\mathsf U_a^{-1}\Psi_\infty(z)c]\\
 &=\langle\Psi_\infty(x)b,
       \overline A_\infty(\mu)\Psi_\infty(z)c\rangle_{\mathsf E}.
\end{aligned}
\]
Both argument sequences belong to $\mathsf H_a$ by
\eqref{eq:full-weighted-feature-norm}; finite partial sums converge there,
so the displayed equality also identifies the infinite coordinate series.
\end{proof}

\begin{lemma}[Weighted feature bounds with explicit conditional moment dependence]
\label{lemma-weighted-feature-envelope-derived}
Assume \eqref{eq:weighted-feature-conditional-input}. Recall
$\chi(u)=[\log(e+u)]^{\kappa_{\rm dec}/(\kappa_{\rm dec}-1)}$ from
\eqref{eq:weighted-series-growth-function}, and define the single fourth-moment envelope
\begin{equation}
\label{eq:conditional-weighted-fourth-envelope}
 C_{\Phi,4}(K):=
 \left[d\sum_{m\ge1}a_m^2L_m^2
 \left\{1+K^{m-1}(4(m-1)\vee1)^{(m-1)/\alpha}\right\}^2\right]^2.
\end{equation}
All constants $C$ below depend only on the fixed feature parameters,
$d$, and $\alpha$, not on $K_{\rm pop},N,t,D$ or any probability budget.
The following assertions hold.
\begin{enumerate}
 \item The full feature map satisfies
 \begin{equation}
 \label{eq:derived-weighted-feature-growth}
 \|\Psi_\infty(x)\|_{\mathrm{HS};\,\mathbb R^d\to\mathsf E}^2
 \le C\exp\{C\chi(\|x\|_{\mathbb R^d})\}.
 \end{equation}
 The feature map is continuous in Hilbert--Schmidt norm and strongly
 measurable.
 \item Uniformly in $t$,
 \begin{equation}
 \label{eq:derived-weighted-fourth-moment}
 \mathbb E_{\mathcal G}\|\Psi_\infty(Y_t^1)\|_{\mathrm{HS};\,\mathbb R^d\to\mathsf E}^4
 \le C_{\Phi,4}(K_{\rm pop})\le C e^{C\chi(K_{\rm pop})}.
 \end{equation}
 \item The Bochner integral
 $\Sigma_t:=\mathbb E_{\mathcal G}[\Psi_\infty(Y_t^1)\Psi_\infty(Y_t^1)^*]$
 is a positive trace-class operator, and
 \begin{equation}
 \label{eq:derived-weighted-covariance-bound}
 \Sigma_t\preceq C_{\Phi,4}(K_{\rm pop})^{1/2}I_{\mathsf E}.
 \end{equation}
\end{enumerate}
\end{lemma}
\begin{proof}
\noindent\textbf{Step 1: continuity and the pointwise bound.}
For $R\ge1$ and $\|x\|_{\mathbb R^d},\|y\|_{\mathbb R^d}\le R$, integrate $\mathrm D^2\phi_m$ on the line
segment from $y$ to $x$. The derivative-growth assumption gives
\[
 \|\mathrm D\phi_m(x)-\mathrm D\phi_m(y)\|_{\mathrm{HS};\,\mathbb R^d\to\mathsf H_m}
 \le\sqrt d\,L_m(1+R^{(m-2)_+})\|x-y\|_{\mathbb R^d}.
\]
After summing squares with weights $a_m^2$, this yields
\[
 \|\Psi_\infty(x)-\Psi_\infty(y)\|_{\mathrm{HS};\,\mathbb R^d\to\mathsf E}^2
 \le d\|x-y\|_{\mathbb R^d}^2\sum_{m\ge1}a_m^2L_m^2(1+R^{(m-2)_+})^2.
\]
The coefficient is finite by squaring \eqref{eq:weighted-series-summand-bound}
with $K=1$, $D=R$ and summing. This proves continuity on every ball. Its image
lies in a separable closed subspace, since $\mathbb R^d$ is separable and
the map is continuous. Strong measurability means being an almost-sure norm limit of measurable
simple functions. A continuous map from the separable space $\mathbb R^d$
has separable image, so this property holds for $\Psi_\infty(Y_t^1)$
without assuming that the whole feature space is separable.
For the pointwise bound, $1+\|x\|_{\mathbb R^d}^{m-1}\le2+(1+\|x\|_{\mathbb R^d})^m$.
Apply \cref{lem:weighted-series-estimate},
\eqref{eq:weighted-series-summand-bound}, with $K=1$ and $D=1+\|x\|_{\mathbb R^d}$:
\begin{equation}
\label{eq:weighted-feature-growth-series}
 \|\Psi_\infty(x)\|_{\mathrm{HS};\,\mathbb R^d\to\mathsf E}^2
 \le d\sum_m[a_mL_m\{2+(1+\|x\|_{\mathbb R^d})^m\}]^2
 \le Ce^{C\chi(\|x\|_{\mathbb R^d})}\sum_m e^{-(c_a/2)m^{\kappa_{\rm dec}}}.
\end{equation}
Here \eqref{eq:weighted-series-scale-comparison} with $c=1$ absorbs
$\chi(2+\|x\|_{\mathbb R^d})$. The finite sum proves
\eqref{eq:derived-weighted-feature-growth}.

\par\noindent\textbf{Step 2: the fourth moment.}
For $m\ge2$, \eqref{eq:weighted-feature-conditional-input} with $r=4(m-1)$
gives
\[
 \|1+\|Y_t^1\|_{\mathbb R^d}^{m-1}\|_{L^4(\mathbb P_{\mathcal G})}
 \le1+(\mathbb E_{\mathcal G}\|Y_t^1\|_{\mathbb R^d}^{4(m-1)})^{1/4}
 \le1+K_{\rm pop}^{m-1}(4(m-1))^{(m-1)/\alpha}.
\]
For $m=1$ the left side is $2$, matching the convention in
\eqref{eq:conditional-weighted-fourth-envelope}. Minkowski's inequality
for a finite sum in $L^2(\mathbb P_{\mathcal G})$ now gives
\[
\begin{aligned}
 \left\|\sum_{m=1}^M a_m^2
       \|\mathrm D\phi_m(Y_t^1)\|_{\mathrm{HS};\,\mathbb R^d\to\mathsf H_m}^2\right\|_{L^2(\mathbb P_{\mathcal G})}
 &\le d\sum_{m\le M}a_m^2L_m^2
       \|1+\|Y_t^1\|_{\mathbb R^d}^{m-1}\|_{L^4(\mathbb P_{\mathcal G})}^2\\
 &\le d\sum_{m\le M}a_m^2L_m^2
 \{1+K_{\rm pop}^{m-1}(4(m-1)\vee1)^{(m-1)/\alpha}\}^2.
\end{aligned}
\]
Squaring and applying monotone convergence gives the explicit connection
between the finite sums and the full fourth moment:
\begin{equation}
\label{eq:weighted-fourth-moment-limit}
\begin{aligned}
 \mathbb E_{\mathcal G}\|\Psi_\infty(Y_t^1)\|_{\mathrm{HS};\,\mathbb R^d\to\mathsf E}^4
 &=\lim_{M\to\infty}\mathbb E_{\mathcal G}
 \left[\sum_{m=1}^M a_m^2
          \|\mathrm D\phi_m(Y_t^1)\|_{\mathrm{HS};\,\mathbb R^d\to\mathsf H_m}^2\right]^2\\
 &\le C_{\Phi,4}(K_{\rm pop}).
\end{aligned}
\end{equation}
For every $m\ge1$, including $m=1$, the braces in
\eqref{eq:conditional-weighted-fourth-envelope} satisfy
\[
 1+K^{m-1}(4(m-1)\vee1)^{(m-1)/\alpha}
 \le 1+\{4^{1/\alpha}(1+K)\}^m m^{m/\alpha}.
\]
Use \eqref{eq:weighted-series-summand-bound} with
$K$ replaced by $4^{1/\alpha}(1+K_{\rm pop})$ and $D=1$.
Squaring each summand, summing, and then squaring the sum in
\eqref{eq:conditional-weighted-fourth-envelope} yields
\[
 C_{\Phi,4}(K_{\rm pop})
 \le\left[Ce^{C\chi(K_{\rm pop})}
           \sum_m e^{-(c_a/2)m^{\kappa_{\rm dec}}}\right]^2
 \le Ce^{C\chi(K_{\rm pop})},
\]
where \eqref{eq:weighted-series-scale-comparison} absorbs the fixed
rescaling. This proves \eqref{eq:derived-weighted-fourth-moment}.
\par\noindent\textbf{Step 3: the covariance.}
For Hilbert--Schmidt maps $A,B:\mathbb R^d\to\mathsf E$, its singular
values $s_j$ give
\[
 \operatorname{tr}(AA^*)=\sum_js_j^2=\|A\|_{\mathrm{HS};\,\mathbb R^d\to\mathsf E}^2,
 \qquad
 \|AA^*\|_{\mathrm{HS};\,\mathsf E\to\mathsf E}=(\sum_js_j^4)^{1/2}\le\sum_js_j^2.
\]
Moreover, $AA^*-BB^*=(A-B)A^*+B(A-B)^*$ gives
\[
 \|AA^*-BB^*\|_{*;\,\mathsf E\to\mathsf E}
 \le(\|A\|_{\mathrm{HS};\,\mathbb R^d\to\mathsf E}+\|B\|_{\mathrm{HS};\,\mathbb R^d\to\mathsf E})\|A-B\|_{\mathrm{HS};\,\mathbb R^d\to\mathsf E}.
\]
Thus $A\mapsto AA^*$ is continuous from Hilbert--Schmidt norm to trace
norm, using \eqref{eq:nuclear-hilbert-schmidt-product}.
Composition with the strongly measurable feature map from Step 1
therefore gives strong measurability in trace norm. For integrability,
positivity of $\Psi_\infty(Y_t^1)\Psi_\infty(Y_t^1)^*$ and
Cauchy--Schwarz give
\[
\begin{aligned}
 \mathbb E_{\mathcal G}
   \|\Psi_\infty(Y_t^1)\Psi_\infty(Y_t^1)^*\|_{*;\,\mathsf E\to\mathsf E}
 &=\mathbb E_{\mathcal G}\|\Psi_\infty(Y_t^1)\|_{\mathrm{HS};\,\mathbb R^d\to\mathsf E}^2\\
 &\le\bigl(\mathbb E_{\mathcal G}
                  \|\Psi_\infty(Y_t^1)\|_{\mathrm{HS};\,\mathbb R^d\to\mathsf E}^4\bigr)^{1/2}
 \le C_{\Phi,4}(K_{\rm pop})^{1/2}<\infty.
\end{aligned}
\]
This establishes Bochner integrability in trace norm.
For $u\in\mathsf E$, Cauchy--Schwarz gives
\[
 \langle u,\Sigma_tu\rangle
 =\mathbb E_{\mathcal G}\|\Psi_\infty(Y_t^1)^*u\|_{\mathbb R^d}^2
 \le\|u\|_{\mathsf E}^2\mathbb E_{\mathcal G}\|\Psi_\infty(Y_t^1)\|_{\mathrm{HS};\,\mathbb R^d\to\mathsf E}^2
 \le C_{\Phi,4}(K_{\rm pop})^{1/2}\|u\|_{\mathsf E}^2.
\]
\end{proof}

\paragraph{The full residual as a feature covariance error.}
Recall the kernel covariance residual from \cref{lemma-kernel-mismatch-bound}:
\begin{equation}
\label{eq:tilde-gamma-full-definition-weighted}
 U_t(z):=[\nabla_W^2F(\mu_t;Y_t^i,z)]_{i=1}^N,
 \qquad
 \widetilde\Gamma_t^N
 =\frac1N\left\|\int U_t(z)U_t(z)^\top
                          (\mu_t-\rho_{Y_t})(dz)\right\|_{*;\,\mathbb R^{dN}\to\mathbb R^{dN}}.
\end{equation}
The matrix $U_t(z)$ has $N$ block rows and $d$ columns. Define
\[
 \Psi_t^{Y,N}(v_1,\ldots,v_N):=\sum_{i=1}^N\Psi_\infty(Y_t^i)v_i,
 \qquad
 \Psi_t^{Y,N}:\mathbb R^{dN}\longrightarrow\mathsf E,
\]
where the domain has the ordinary Euclidean inner product. Recall $\Sigma_t$
from \cref{lemma-weighted-feature-envelope-derived} and define
\[
 \widehat\Sigma_t^N:=\frac1N\sum_{i=1}^N
       \Psi_\infty(Y_t^i)\Psi_\infty(Y_t^i)^*
       =\frac1N\Psi_t^{Y,N}(\Psi_t^{Y,N})^*.
\]

\begin{lemma}[Full covariance representation and nuclear-norm bound]
\label{lem:tilde-gamma-weighted-representation}
The full residual satisfies
\begin{align}
 \widetilde\Gamma_t^N
 &=\frac1N\left\|(\Psi_t^{Y,N})^*\overline A_\infty(\mu_t)
       (\Sigma_t-\widehat\Sigma_t^N)
       \overline A_\infty(\mu_t)\Psi_t^{Y,N}\right\|_{*;\,\mathbb R^{dN}\to\mathbb R^{dN}},
 \label{eq:tilde-gamma-weighted-matrix-identity}\\
 \widetilde\Gamma_t^N
 &\le\mathrm L_D^2\operatorname{tr}(\widehat\Sigma_t^N)
                  \|\widehat\Sigma_t^N-\Sigma_t\|_{\mathrm{HS};\,\mathsf E\to\mathsf E}.
 \label{eq:tilde-gamma-deterministic-reduction}
\end{align}
Consequently,
\begin{equation}
\label{eq:tilde-gamma-full-expectation-bound}
 \mathbb E_{\mathcal G}\widetilde\Gamma_t^N
 \le\frac{\mathrm L_D^2 C_{\Phi,4}(K_{\rm pop})}{\sqrt N}.
\end{equation}
\end{lemma}
\begin{proof}
\noindent\textbf{Step 1: factor and bound the covariance residual.}
By \eqref{eq:full-weighted-hessian-factorization},
\[
 U_t(z)=(\Psi_t^{Y,N})^*\overline A_\infty(\mu_t)\Psi_\infty(z).
\]
Since $\overline A_\infty(\mu_t)$ is self-adjoint, the product is
\[
 U_t(z)U_t(z)^\top
 =(\Psi_t^{Y,N})^*\overline A_\infty(\mu_t)
 [\Psi_\infty(z)\Psi_\infty(z)^*]
 \overline A_\infty(\mu_t)\Psi_t^{Y,N}.
\]
Only the bracketed factor depends on $z$. Integration of this factor
$\Psi_\infty(z)\Psi_\infty(z)^*$ against $\mu_t$ gives $\Sigma_t$,
and integration against $\rho_{Y_t}=N^{-1}\sum_i\delta_{Y_t^i}$ gives
$\widehat\Sigma_t^N$. Subtracting proves
\eqref{eq:tilde-gamma-weighted-matrix-identity}.
Apply the Hilbert--Schmidt product bound
\eqref{eq:nuclear-hilbert-schmidt-product} to the two factors
$(\Psi_t^{Y,N})^*\overline A_\infty(\mu_t)$ and
$(\Sigma_t-\widehat\Sigma_t^N)\overline A_\infty(\mu_t)\Psi_t^{Y,N}$.
Then use submultiplicativity of the Hilbert--Schmidt norm with the
operator norm, also after taking adjoints:
\[
\begin{aligned}
 \widetilde\Gamma_t^N
 &\le\frac1N\|(\Psi_t^{Y,N})^*\overline A_\infty(\mu_t)\|_{\mathrm{HS};\,\mathsf E\to\mathbb R^{dN}}
 \|(\Sigma_t-\widehat\Sigma_t^N)\overline A_\infty(\mu_t)
                              \Psi_t^{Y,N}\|_{\mathrm{HS};\,\mathbb R^{dN}\to\mathsf E}\\
 &\le\frac1N\|\Psi_t^{Y,N}\|_{\mathrm{HS};\,\mathbb R^{dN}\to\mathsf E}^2
       \|\overline A_\infty(\mu_t)\|_{\mathsf E\to\mathsf E}^2
       \|\Sigma_t-\widehat\Sigma_t^N\|_{\mathsf E\to\mathsf E}\\
 &\le\mathrm L_D^2\operatorname{tr}(\widehat\Sigma_t^N)
       \|\Sigma_t-\widehat\Sigma_t^N\|_{\mathrm{HS};\,\mathsf E\to\mathsf E}.
\end{aligned}
\]
Here $N^{-1}\|\Psi_t^{Y,N}\|_{\mathrm{HS};\,\mathbb R^{dN}\to\mathsf E}^2
=N^{-1}\sum_i\|\Psi_\infty(Y_t^i)\|_{\mathrm{HS};\,\mathbb R^d\to\mathsf E}^2
=\operatorname{tr}(\widehat\Sigma_t^N)$.

\par\noindent\textbf{Step 2: bound its conditional expectation.}
For the expectation bound put
$W_i=\Psi_\infty(Y_t^i)\Psi_\infty(Y_t^i)^*$ and
$Z_i=\|\Psi_\infty(Y_t^i)\|_{\mathrm{HS};\,\mathbb R^d\to\mathsf E}^2$. Jensen's inequality gives
\[
 \mathbb E_{\mathcal G}\left(\operatorname{tr}\widehat\Sigma_t^N\right)^2
 =\mathbb E_{\mathcal G}(N^{-1}\sum_iZ_i)^2
 \le N^{-1}\sum_i\mathbb E_{\mathcal G}Z_i^2
 \le C_{\Phi,4}(K_{\rm pop}).
\]
In the Hilbert space $\mathcal S_2(\mathsf E)$, conditional independence
and $\mathbb E_{\mathcal G}(W_i-\Sigma_t)=0$ imply, for $i\ne j$,
\[
 \mathbb E_{\mathcal G}\langle W_i-\Sigma_t,W_j-\Sigma_t\rangle_{\mathrm{HS};\,\mathsf E\to\mathsf E}
 =\langle\mathbb E_{\mathcal G}(W_i-\Sigma_t),
          \mathbb E_{\mathcal G}(W_j-\Sigma_t)\rangle_{\mathrm{HS};\,\mathsf E\to\mathsf E}=0.
\]
In particular, the singular-value calculation in
\cref{lemma-weighted-feature-envelope-derived} gives
$\|W_1\|_{\mathrm{HS};\,\mathsf E\to\mathsf E}^2\le(\operatorname{tr}W_1)^2=Z_1^2$,
with an inequality rather than equality in general. Expanding gives
\begin{equation}
\label{eq:weighted-covariance-mean-square}
\begin{aligned}
 \mathbb E_{\mathcal G}\|\widehat\Sigma_t^N-\Sigma_t\|_{\mathrm{HS};\,\mathsf E\to\mathsf E}^2
 &=N^{-2}\sum_{i,j}\mathbb E_{\mathcal G}
        \langle W_i-\Sigma_t,W_j-\Sigma_t\rangle_{\mathrm{HS};\,\mathsf E\to\mathsf E}\\
 &=N^{-2}\sum_i\mathbb E_{\mathcal G}\|W_i-\Sigma_t\|_{\mathrm{HS};\,\mathsf E\to\mathsf E}^2\\
 &=N^{-1}(\mathbb E_{\mathcal G}\|W_1\|_{\mathrm{HS};\,\mathsf E\to\mathsf E}^2-\|\Sigma_t\|_{\mathrm{HS};\,\mathsf E\to\mathsf E}^2)\\
 &\le N^{-1}\mathbb E_{\mathcal G}Z_1^2
 \le C_{\Phi,4}(K_{\rm pop})/N.
\end{aligned}
\end{equation}
Conditional Cauchy--Schwarz in \eqref{eq:tilde-gamma-deterministic-reduction}
now gives
\[
\begin{aligned}
 \mathbb E_{\mathcal G}\widetilde\Gamma_t^N
 &\le\mathrm L_D^2
 [\mathbb E_{\mathcal G}(\operatorname{tr}\widehat\Sigma_t^N)^2]^{1/2}
 [\mathbb E_{\mathcal G}\|\widehat\Sigma_t^N-\Sigma_t\|_{\mathrm{HS};\,\mathsf E\to\mathsf E}^2]^{1/2}\\
 &\le\mathrm L_D^2 C_{\Phi,4}(K_{\rm pop})^{1/2}
          [C_{\Phi,4}(K_{\rm pop})/N]^{1/2}.
\end{aligned}
\]
The two factors use the same sample; their independence is not needed.
\end{proof}

\begin{lemma}[A deterministic envelope at the sampled particles]
\label{lem:weighted-feature-max-event}
Fix a time trace $t$, $N\ge2$, and $\delta\in(0,1)$. Define
\[
 R_{N,\delta}:=K_{\rm pop}
   [e(\alpha\vee1)\{1+\log(8N/\delta)\}]^{1/\alpha},
\]
\begin{equation}
\label{eq:weighted-feature-max-tau-N-delta}
 \tau_{N,\delta}:=1\vee
   \left[d\sum_{m\ge1}a_m^2L_m^2(1+R_{N,\delta}^{m-1})^2\right].
\end{equation}
Then
\begin{equation}
\label{eq:weighted-feature-max-event}
 \mathbb P_{\mathcal G}\left(
   \max_{i\le N}\|\Psi_\infty(Y_t^i)\|_{\mathrm{HS};\,\mathbb R^d\to\mathsf E}^2>\tau_{N,\delta}
                         \right)\le\delta/8,
\end{equation}
and the individual-particle probability is at most $\delta/(8N)$.
Moreover $\tau_{N,\delta}\le C\exp\{C\chi(R_{N,\delta})\}$.
If $\Xi\ge N/\delta$ and $K_{\rm pop}\le C_K(\log\Xi)^{s_K}$ for fixed
$C_K,s_K$, then, for every $\varepsilon>0$,
\begin{equation}
\label{eq:weighted-envelope-subpolynomial}
 \tau_{N,\delta}^2\sqrt{\log(8/\delta)}
       +C_{\Phi,4}(K_{\rm pop})\le C_\varepsilon \Xi^\varepsilon.
\end{equation}
\end{lemma}
\begin{proof}
Use the Markov calculation of \cref{lemma-beta-tail-bound},
\eqref{eq:moment-tail-markov-substitution}, with threshold $R_{N,\delta}$.
In the moment input \eqref{eq:weighted-feature-conditional-input}, take
$r=(\alpha\vee1)\{1+\log(8N/\delta)\}\ge1$. Then
\[
 \mathbb P_{\mathcal G}(\|Y_t^i\|_{\mathbb R^d}>R_{N,\delta})
 \le\left(\frac{K_{\rm pop}r^{1/\alpha}}
                   {K_{\rm pop}(er)^{1/\alpha}}\right)^r
 =e^{-r/\alpha}\le\frac{\delta}{8N}.
\]
By \eqref{eq:full-weighted-feature-norm} and
\eqref{eq:weighted-feature-max-tau-N-delta},
\[
 \{\|\Psi_\infty(Y_t^i)\|_{\mathrm{HS};\,\mathbb R^d\to\mathsf E}^2>\tau_{N,\delta}\}
 \subseteq\{\|Y_t^i\|_{\mathbb R^d}>R_{N,\delta}\}.
\]
The union over $N$ particles costs $N\delta/(8N)=\delta/8$, proving
\eqref{eq:weighted-feature-max-event}.
The calculation \eqref{eq:weighted-feature-growth-series} in
\cref{lemma-weighted-feature-envelope-derived}, with $\|x\|_{\mathbb R^d}$ replaced by
$R_{N,\delta}$, proves the growth bound.
Since $\Xi\ge N/\delta\ge2$,
$\log(8N/\delta)\le\log8+\log\Xi\le4\log\Xi$. Hence
\[
 R_{N,\delta}\le C(\log\Xi)^{s_K+1/\alpha},\qquad
 \log(e+R_{N,\delta})\le C+C\log\log\Xi.
\]
For the whole left-hand side of \eqref{eq:weighted-envelope-subpolynomial},
these bounds and \eqref{eq:derived-weighted-fourth-moment} give
\[
\begin{aligned}
 &\log\{\tau_{N,\delta}^2\sqrt{\log(8/\delta)}
                      +C_{\Phi,4}(K_{\rm pop})\}\\
 &\quad\le C+C[1+\log\log(e+\Xi)]^{\kappa_{\rm dec}/(\kappa_{\rm dec}-1)}
 \le C_\varepsilon+\varepsilon\log\Xi.
\end{aligned}
\]
The last inequality holds because any fixed power of $\log\log(e+\Xi)$
is $o(\log\Xi)$; a fixed constant covers bounded $\Xi\ge2$.
Exponentiation proves the claim. The constant may depend on $C_K,s_K$
and $\varepsilon$, but not on $\Xi$.
\end{proof}

\begin{lemma}[Full weighted covariance concentration in Hilbert--Schmidt norm]
\label{lem:tilde-gamma-weighted-covariance-input}
Fix a time trace $t$, $N\ge2$, and $\delta\in(0,1)$. Use the threshold
$\tau_{N,\delta}$ in \eqref{eq:weighted-feature-max-tau-N-delta}.
There is an event of conditional probability at least $1-\delta$ on which
\begin{align}
 \max_{i\le N}\|\Psi_\infty(Y_t^i)\|_{\mathrm{HS};\,\mathbb R^d\to\mathsf E}^2&\le\tau_{N,\delta},
 \label{eq:full-covariance-sample-envelope}\\
 \|\widehat\Sigma_t^N-\Sigma_t\|_{\mathrm{HS};\,\mathsf E\to\mathsf E}
 &\le\frac{C_{\Phi,4}(K_{\rm pop})^{1/2}}{\sqrt N}
          (1+\sqrt{\delta/8})
       +\tau_{N,\delta}\sqrt{\frac{2\log(8/\delta)}N}.
 \label{eq:tilde-gamma-covariance-good-event}
\end{align}
\end{lemma}
\begin{proof}
The repeated random objects needed in this proof are
\[
 W_i:=\Psi_\infty(Y_t^i)\Psi_\infty(Y_t^i)^*,\quad
 Z_i:=\|\Psi_\infty(Y_t^i)\|_{\mathrm{HS};\,\mathbb R^d\to\mathsf E}^2,\quad
 W_i^{(\tau)}:=W_i\mathbf1_{\{Z_i\le\tau_{N,\delta}\}}.
\]
The conditional mean of $W_i^{(\tau)}$ is
$\Sigma_t^{(\tau)}:=\mathbb E_{\mathcal G}W_1^{(\tau)}$.
The sample mean of $W_i^{(\tau)}$ is
$\widehat\Sigma_t^{N,(\tau)}:=N^{-1}\sum_iW_i^{(\tau)}$.
Conditional on $\mathcal G_{\rm pop}$, these clipped variables remain
independent. We apply concentration before intersecting with
$\{\max_iZ_i\le\tau_{N,\delta}\}$.

\par\noindent\textbf{Step 1: mean square and bounded differences.}
The singular-value inequality in
\cref{lemma-weighted-feature-envelope-derived} gives
$\|W_i^{(\tau)}\|_{\mathrm{HS};\,\mathsf E\to\mathsf E}\le Z_i\mathbf1_{\{Z_i\le\tau_{N,\delta}\}}
\le\tau_{N,\delta}$. These operators are conditionally i.i.d. and
square integrable. Apply the centered-sum calculation
\eqref{eq:weighted-covariance-mean-square} of
\cref{lem:tilde-gamma-weighted-representation}, replacing
$W_i,\Sigma_t$ by $W_i^{(\tau)},\Sigma_t^{(\tau)}$:
\[
 \mathbb E_{\mathcal G}
 \|\widehat\Sigma_t^{N,(\tau)}-\Sigma_t^{(\tau)}\|_{\mathrm{HS};\,\mathsf E\to\mathsf E}^2
 \le N^{-1}\mathbb E_{\mathcal G}\|W_1^{(\tau)}\|_{\mathrm{HS};\,\mathsf E\to\mathsf E}^2
 \le N^{-1}\mathbb E_{\mathcal G}Z_1^2
 \le C_{\Phi,4}(K_{\rm pop})/N.
\]
Cauchy--Schwarz therefore bounds the mean norm by
$\sqrt{C_{\Phi,4}(K_{\rm pop})/N}$.
If sample $i$ is replaced, denote the new sample mean by
$\widehat\Sigma_t^{N,(\tau)\prime}$ and its changed summand by
$W_i^{(\tau)\prime}$. The threshold and population mean are fixed, so
\[
\begin{aligned}
 &\left|\|\widehat\Sigma_t^{N,(\tau)}-\Sigma_t^{(\tau)}\|_{\mathrm{HS};\,\mathsf E\to\mathsf E}
 -\|\widehat\Sigma_t^{N,(\tau)\prime}-\Sigma_t^{(\tau)}\|_{\mathrm{HS};\,\mathsf E\to\mathsf E}\right|\\
 &\quad\le\|\widehat\Sigma_t^{N,(\tau)}-
                        \widehat\Sigma_t^{N,(\tau)\prime}\|_{\mathrm{HS};\,\mathsf E\to\mathsf E}
 =N^{-1}\|W_i^{(\tau)}-W_i^{(\tau)\prime}\|_{\mathrm{HS};\,\mathsf E\to\mathsf E}
 \le 2\tau_{N,\delta}/N.
\end{aligned}
\]
In \cref{lemma:mcdiarmid}, the sum of squared differences is
$N(2\tau_{N,\delta}/N)^2=4\tau_{N,\delta}^2/N$. Substituting
$s=\tau_{N,\delta}\sqrt{2\log(8/\delta)/N}$ gives failure probability
\[
 \exp\left(-\frac{2s^2}{4\tau_{N,\delta}^2/N}\right)
 =\exp(-\log(8/\delta))=\delta/8.
\]
Thus, except on that event,
\[
 \|\widehat\Sigma_t^{N,(\tau)}-\Sigma_t^{(\tau)}\|_{\mathrm{HS};\,\mathsf E\to\mathsf E}
 \le\sqrt{C_{\Phi,4}(K_{\rm pop})/N}
       +\tau_{N,\delta}\sqrt{2\log(8/\delta)/N}.
\]

\par\noindent\textbf{Step 2: restore the conditional mean and the sample.}
The mean difference is
$\Sigma_t-\Sigma_t^{(\tau)}=\mathbb E_{\mathcal G}[W_1\mathbf1_{\{Z_1>\tau_{N,\delta}\}}]$.
The individual tail bound in \cref{lem:weighted-feature-max-event}
and Cauchy--Schwarz give
\[
\begin{aligned}
 \|\Sigma_t-\Sigma_t^{(\tau)}\|_{\mathrm{HS};\,\mathsf E\to\mathsf E}
 &\le\mathbb E_{\mathcal G}
       [\|W_1\|_{\mathrm{HS};\,\mathsf E\to\mathsf E}\mathbf1_{\{Z_1>\tau_{N,\delta}\}}]\\
 &\le(\mathbb E_{\mathcal G}Z_1^2)^{1/2}
       \mathbb P_{\mathcal G}(Z_1>\tau_{N,\delta})^{1/2}
 \le C_{\Phi,4}(K_{\rm pop})^{1/2}\sqrt{\delta/(8N)}.
\end{aligned}
\]
The bound \eqref{eq:weighted-feature-max-event} in
\cref{lem:weighted-feature-max-event} costs another $\delta/8$ to ensure
$\max_iZ_i\le\tau_{N,\delta}$ and hence
$\widehat\Sigma_t^N=\widehat\Sigma_t^{N,(\tau)}$. On the intersection,
\[
\begin{aligned}
 \|\widehat\Sigma_t^N-\Sigma_t\|_{\mathrm{HS};\,\mathsf E\to\mathsf E}
 &\le\|\widehat\Sigma_t^{N,(\tau)}-\Sigma_t^{(\tau)}\|_{\mathrm{HS};\,\mathsf E\to\mathsf E}
          +\|\Sigma_t^{(\tau)}-\Sigma_t\|_{\mathrm{HS};\,\mathsf E\to\mathsf E}\\
 &\le\frac{C_{\Phi,4}(K_{\rm pop})^{1/2}}{\sqrt N}(1+\sqrt{\delta/8})
       +\tau_{N,\delta}\sqrt{\frac{2\log(8/\delta)}N}.
\end{aligned}
\]
The population bias bound is deterministic after conditioning and costs
no further probability. Thus the failure budget is
$\delta/8+\delta/8\le\delta$, proving
\eqref{eq:tilde-gamma-covariance-good-event}.
\end{proof}

\begin{corollary}[The full jump residual at a fixed time]
\label{cor:tilde-gamma-full-weighted-bound}
For every fixed time trace $t$, $N\ge2$, and $\delta\in(0,1)$, with
conditional probability at least $1-\delta$,
\begin{equation}
\label{eq:tilde-gamma-full-high-prob-bound}
 \widetilde\Gamma_t^N
 \le\frac{\mathrm L_D^2\tau_{N,\delta}}{\sqrt N}
       \left[2C_{\Phi,4}(K_{\rm pop})^{1/2}
                 +\tau_{N,\delta}\sqrt{2\log(8/\delta)}\right].
\end{equation}
\end{corollary}
\begin{proof}
On the joint event of \cref{lem:tilde-gamma-weighted-covariance-input},
\[
 \operatorname{tr}(\widehat\Sigma_t^N)
 =N^{-1}\sum_i\|\Psi_\infty(Y_t^i)\|_{\mathrm{HS};\,\mathbb R^d\to\mathsf E}^2\le\tau_{N,\delta}.
\]
Substitute this and \eqref{eq:tilde-gamma-covariance-good-event} into
\cref{lem:tilde-gamma-weighted-representation},
\eqref{eq:tilde-gamma-deterministic-reduction}:
\[
 \widetilde\Gamma_t^N\le\mathrm L_D^2\tau_{N,\delta}
 \left[\frac{C_{\Phi,4}(K_{\rm pop})^{1/2}}{\sqrt N}(1+\sqrt{\delta/8})
       +\tau_{N,\delta}\sqrt{\frac{2\log(8/\delta)}N}\right].
\]
Use $1+\sqrt{\delta/8}\le2$ to obtain
\eqref{eq:tilde-gamma-full-high-prob-bound} on the same event.
\end{proof}

\begin{corollary}[Union bound over prescribed jump times]
\label{cor:tilde-gamma-jump-time-union-bound}
Let $0<t_1<\cdots<t_{\NofJ}\le T$ be deterministic, and set
$q=p/(\NofJ+1)$ for $p\in(0,1)$. Define
\begin{equation}
\label{eq:phase1-HS-tilde-delta-envelope}
 \widetilde\delta_{N,T,p}^{\rm HS}
 :=\frac{\mathrm L_D^2\tau_{N,q}}{\sqrt N}
       \left[2C_{\Phi,4}(K_{\rm pop})^{1/2}
                    +\tau_{N,q}\sqrt{2\log(8/q)}\right].
\end{equation}
Then, with conditional probability at least $1-p$,
\begin{equation}
\label{eq:tilde-gamma-jump-time-union-bound}
 \max_{1\le\ell\le\NofJ}\widetilde\Gamma_{t_\ell-}^N
 \le\widetilde\delta_{N,T,p}^{\rm HS}.
\end{equation}
The maximum is zero when $\NofJ=0$. The parameters $K_{\rm pop}$ and
$\mathrm L_D$ retain their horizon and population-confidence dependence.
If $K_{\rm pop},D$ grow at most as fixed powers of $\log\Xi$,
$\Xi\ge N(\NofJ+1)/p$, and the structural parameters are fixed, then
for every $\varepsilon>0$,
\begin{equation}
\label{eq:tilde-gamma-jump-time-final-rate}
 \widetilde\delta_{N,T,p}^{\rm HS}
 \le C_\varepsilon N^{-1/2}\Xi^\varepsilon.
\end{equation}
In particular the residual has rate $N^{-1/2+o(1)}$ when $\Xi$ is polynomial
in $N$; its square root in the jump comparison has rate $N^{-1/4+o(1)}$.
\end{corollary}
\begin{proof}
\noindent\textbf{Step 1: combine the prescribed jump times.}
Apply \cref{cor:tilde-gamma-full-weighted-bound} at each deterministic
pre-jump trace with $\delta=q=p/(\NofJ+1)$ in
\eqref{eq:tilde-gamma-full-high-prob-bound}. The union costs
$\NofJ q=p\NofJ/(\NofJ+1)\le p$; different times need not be independent.

\par\noindent\textbf{Step 2: simplify the deterministic coefficient.}
For the rate, \eqref{eq:local-envelope-growth} bounds
$\mathrm L_D\le Ce^{C\chi(K_{\rm pop}+D)}$.
For polylogarithmic $K_{\rm pop},D$, this is at most
$C_\eta \Xi^\eta$ for every $\eta>0$, by the logarithm calculation proving
\eqref{eq:weighted-envelope-subpolynomial} in
\cref{lem:weighted-feature-max-event}. Apply
\eqref{eq:weighted-envelope-subpolynomial} with the same $\eta$.
Since $\log(8/q)>1$, it implies
\[
 \tau_{N,q}\le C_\eta \Xi^{\eta/2},\quad
 C_{\Phi,4}(K_{\rm pop})^{1/2}\le C_\eta \Xi^{\eta/2},\quad
 \tau_{N,q}^2\sqrt{\log(8/q)}\le C_\eta \Xi^\eta.
\]
The bracket after multiplication by $\tau_{N,q}$ is therefore at most
$C_\eta \Xi^\eta$. Together with $\mathrm L_D^2\le C_\eta \Xi^{2\eta}$,
this gives \eqref{eq:tilde-gamma-jump-time-final-rate} by taking
$\eta=\varepsilon/3$.

\par\noindent\textbf{Step 3: average the conditional failure probability.}
Finally, since $E_{\rm pop}\in\mathcal G_{\rm pop}$,
\[
\begin{aligned}
 &\mathbb P\left(E_{\rm pop}\cap
   \{\max_\ell\widetilde\Gamma_{t_\ell-}^N
                  >\widetilde\delta_{N,T,p}^{\rm HS}\}\right)\\
 &\quad=\mathbb E\left[\mathbf1_{E_{\rm pop}}
       \mathbb P_{\mathcal G}\left(
          \max_\ell\widetilde\Gamma_{t_\ell-}^N
                     >\widetilde\delta_{N,T,p}^{\rm HS}\right)\right]
 \le p.
\end{aligned}
\]
Adding $\mathbb P(E_{\rm pop}^c)$ gives the unconditional failure budget
when the population event is also required.
\end{proof}

\begin{remark}[Why no feature cutoff or nuclear-norm rank loss is needed]
\label{remark:tilde-gamma-no-rank-loss}
The finite $dN\times dN$ residual factors through the full Hilbert feature
space $\mathsf E$. The independent random operators are clipped by their
norm before scalar bounded differences is applied in $\mathcal S_2(\mathsf E)$.
Neither argument uses the dimension of a feature space. The nuclear norm
is bounded by the empirical feature trace times the Hilbert--Schmidt
covariance error on one event. This proves the bound for
$\widetilde\Gamma_t^N$ directly.
\end{remark}
